% ==================== 摘要专用页 ====================
\begin{center}
  {\zihao{3}\heiti 基于YOLOv8与极值理论的集装箱破损智能检测模型}
\end{center}

\vspace{1.5em}

\noindent{\heiti\zihao{-4} 摘要：}

\zihao{-4}
针对集装箱外表面破损的智能检测问题，本文构建了“检测—分类—评估”一体化的计算机视觉建模框架，依次解决破损有无判别（问题一）、破损定位与类别识别（问题二）以及模型的多维度评估（问题三）三项任务。

\vspace{0.35em}
对于问题一，考虑到训练集全部为含破损图像、缺乏无破损负样本，直接训练二分类网络难以刻画“无破损”类别。本文提出基于极值理论（Extreme Value Theory, EVT）的开放集判别模型：利用检测网络提取每幅图像的最大检测置信度作为统计量，依据 Fisher--Tippett--Gnedenko 定理以二参数 Weibull 分布对其极值分布进行建模，并以分布的累积概率定义缺陷概率，结合伪负样本校准确定最优阈值。最终模型拟合参数为 $k=3.064$、$\lambda=0.701$，正负样本最大置信度均值（0.628 对 0.077）分离显著，近似 AUC 达 0.972，等价置信度 0.34 处召回率 85.2\%、误报率仅 5.3\%；最终模型在测试集上判别出 352 幅破损图像与 61 幅无破损图像，判决阈值与检测输出保持统计一致性。

\vspace{0.35em}
对于问题二，本文基于 YOLOv8 构建缺陷定位与分类模型。针对 Hole（破洞）类样本占比仅 11.8\% 的类别不平衡问题，采用针对性数据增强策略（含 Copy-Paste 等）与类别均衡设计；针对训练集无负样本导致背景误检的问题，从训练图像无标注区域裁剪 600 幅负样本参与训练。三组对照实验表明：骨干网络升级（YOLOv8n$\to$YOLOv8s）与数据增强策略使验证集 mAP@0.5:0.95 由 0.190 提升至 0.212；引入负样本训练的最终模型在验证集上达到 mAP@0.5 = 0.405、mAP@0.5:0.95 = 0.205，逐类 mAP@0.5:0.95 为 Dent 0.295、Hole 0.190、Rusty 0.129，其中 Rusty 类召回率仅 0.282 是主要短板；负样本比例实验验证 600 幅为分离度甜点，检测结果按竞赛要求输出至 test\_result.csv。

\vspace{0.35em}
对于问题三，本文从精度、分类、消融、鲁棒性、效率与错误分析六个维度系统评估模型：检测精度上，独立复测与训练最优指标一致，Rusty 类弱纹理缺陷是主要瓶颈；消融实验上，增益主要来自骨干升级，Copy-Paste 贡献可忽略，负样本训练以约 0.007 的 mAP 代价改善分离度；鲁棒性上，模型对亮度（$\pm30$）与 $3\times3$ 高斯模糊扰动稳健（mAP@0.5:0.95 衰减不超过 0.014），对强高斯噪声敏感（$\sigma=50$ 时降至 0.004）；效率上，最终模型参数规模约 11.1 M、计算量 28.4 GFLOPs，在 RTX 4070 Laptop GPU 上单幅 640$\times$640 图像端到端推理 8.87 ms（约 113 FPS），满足实时巡检需求；错误分析上，小目标漏检与类间混淆是主要误差来源；P2 头使 Hole 类 mAP@0.5:0.95 提升 0.021、TTA 使整体指标提升至 0.215，为后续优化指明方向。

\vspace{1em}

\noindent{\heiti\zihao{-4} 关键词：}
{\zihao{-4} 集装箱破损检测；YOLOv8；极值理论；Weibull 分布；Copy-Paste 数据增强；多维度评估}

\clearpage
